\section{Lair of \npcMidnightMother{}}
\label{sec:middnightMother}

The source of all wrongness affecting the Tower, murderers responsible for half of the deaths, lurk in here. After the long search, the PCs can finally free \locationUrithiru{} from it's influence.




\subsection{Finding the Entrance}

To find the entrance to the Lair, the PCs will need to investigate murders committed by the Midnight Essence (see \fullref{sec:mysteriousViolence}) and/or search for it in the labyrinth of \locationUrithiru{} tunnels (see \fullref{sec:darkTunnels}.) 

When the party finds the entrance, they may immediately go inside, or they might go back for reinforcements. If they enter, proceed to \fullref{ssec:middnightMotherEnter}. 

If they decide to gather more allies, they will need to find the entrance again while descending into \fullref{sec:darkTunnels}. However, this time their skill checks to find it gain an advantage.




\subsection{No Backing Down}
\label{ssec:middnightMotherNoBacking}

At some point, after \fullref{ssec:mysteriousViolenceBarFight}, \npcShallan{} will be able to find this location on her own. When she does, she will send for \npcAdolin{} who will arrive with \npcRenarin{} and several bridgemen. 

At this point, the party can join them. If they do, read the following:

\quoteBox{
	"I brought the best men I could find", you hear Brightlord \npcAdolin{} saying. "But I feel I should bring an entire army instead. You sure you want to o this now?"
}

At first, \npcShallan{} will be hesitant weather she should enter inside. If a PC would attempt to persuade her to wait (perhaps to gather more allies), she will agree without any rolls required. 

If that happens, count the scene as an additional success towards their Endeavor to find \npcMidnightMother{}. That means they might need to do more searching in the future.

If the PCs agree to follow \npcShallan{} and \npcAdolin{} on the spot, proceed to \fullref{ssec:middnightMotherEnter}.

It is also possible for the PCs to bail \npcShallan{} and let her enter without their help. When this happen, allow their spren to warn them of missing the glory, spoils and recognition. If the PCs persist, proceed to \fullref{ssec:middnightMotherIgnored}.


\spoilerBox{
	The same scene is referenced in \bookOB{}, Chapter 29. "No Backing Down".}


\subsection{Descending into the Lair}
\label{ssec:middnightMotherEnter}

When the party descends into the darkness, read the following:

\quoteBox{
	You walk a seemingly endless stairs deep into the darkness of the Tower. Finally, you reach the bottom. Someone takes out a sphere to illuminate the room. In its light, you see ancient mosaics depicting all ten of the Radiant Orders. \\ \\
	But you cannot admire the masterpiece right now... Not with the overwhelming feeling of wrongness emanating from the hall before you.
}

No matter how the entrance was found, even if the party decides to descend without any allies, somehow \npcVeil{} will find out about that and join them. She will switch to \npcShallan{} once the party reaches the bottom of the stairs.

Allow the party to explore the surroundings and prepare for the fight they can sense awaiting for them.

\reasonBox{
	Make sure \npcShallan{} is present during the fight with the Unmade. You can use her to confront \npcMidnightMother{} as a fail-safe if no PC steps up to do it. In other case, the fight would never end...
}

\subsection{\npcMidnightMotherFull{}}
\label{ssec:middnightMotherFight}

When the PCs reach this location, read the following:

\quoteBox{
	At last, the light of your spheres illuminates something before you. A large mass, reflecting light like a glistening tar. \\ \\
	"We shouldn't have come here", you hear someone saying. "We can't fight this. Stormfather..."
}

The party is facing \npcMidnightMotherFull{}. They must defeat her to restore order to the Tower.

\corruptedBox{
	Seeing \npcMidnightMotherFull{} can expedite an Enlightened Edgedancer \seeCorruptedRules{}.
}

\reasonBox{
	This is the last time this adventure will feature this Unmade. Make sure you expedited Enlightened Edgedancers in the party three times in total at this point.
}


\subsubsection{Endless Onslaught}

When the party faces \npcMidnightMother{}, read the following:

\quoteBox{
	Everywhere around, you see a tar-like darkness seeping through the walls and dripping from the ceiling. When the substance hits the ground, it forms a human-shaped monstrosity. Each of them looks like one of you, just... wrong.
}

For each character in the scene, spawn \enemyMidnightAssassin{} fighting that character. At beginning of each round of the combat, spawn additional \enemyMidnightAssassin{} attacking the PCs.

The creatures will fight until the confrontation with \npcMidnightMother{} is completed - see \fullref{ssec:middnightMotherConfront}.

\reasonBox{
	If your players enjoy the combat and want a higher challenge, you can replace two of the \enemyMidnightAssassin{}s facing them with a \enemyMidnightTwin{}.
}

\subsubsection{Fighting \npcMidnightMotherTitle{}}
\label{ssec:middnightMotherConfront}

To stop the endless swarm of the Midnight creatures, one of the characters will need to face \npcMidnightMother{} in a confrontation. If none of the PCs approaches her, \npcShallan{} will do it at the end of the 2nd round of combat.

If a PC chooses to face her, you can run the confrontation as a modified Endeavor. To succeed, the PC must accumulate 4 successes. There is no failing threshold, but each failure means 1 more round of combat the rest of the party must face the enemy.

During each round of combat, the PC confronting the Unmade must make a \skillCheck{25}{\skillIntimiation{}} or a \skillCheck{15}{Illumination}. If the test is passed, increase the total number of successes accumulated for the Endeavor by 1. When the 4th success is gathered, the Unmade is defeated and flees from \locationUrithiru{}.

If the PC feels they should be allies with the Unmade instead of fighting it, allow the PC to make a \skillCheck{25}{\skillPersuasion{}} instead. During this attempt, the PC tries to convince \npcMidnightMother{} she is in a loosing position and must flee to save herself from those Radiants who are now living in the Tower.

The PC who is confronting \npcMidnightMother{} is excluded from the ongoing combat - they can't attack or be attacked by the creatures. However, other PCs can use their actions to support them, eg. by using Aid action.

If no PC steps up to fight the Unmade, assume \npcShallan{} will need 6 rounds to defeat her.


\subsubsection{Goal}

The PC who confronts the Unmade must succeed on 4 skill tests using Illumination, \skillIntimiation{} or \skillPersuasion{}.

The other PCs must survive until the confrontation with \npcMidnightMother{} is finished, at which point all of her creatures will flee from the fight.

\subsubsection{Off-screen Fight}

All allies that are present in the scene are fighting off-screen. Don't bother with rolling for each attack they trade with a \enemyMidnightAssassin{} they are fighting with.

At the end of the fight, assume all of them survived. Some of them could have earned an injury, but they are never in a lethal danger.

\subsubsection{Opportunities and Complications}

During this fight, you can spend Opportunities and Complications in a following way:\\

\oppCompTable{
	\oppCompRow{\iconOpp[20px]{}}{You call one of your allies to help you - see \fullref{ssec:middnightMotherAllies}. If you do, that ally can't be called again in this or following round.}
	\oppCompRow{\iconOpp[20px]{}}{You help one of your allies - ignore the \iconComp{} effect that prevents them from answering a call for help (another \iconOpp{} must be spent to call them).}
	\oppCompRow{\iconComp[20px]{}}{One of the allies is having troubles and can't be called for help during this or the next round of combat.}
	\oppCompRow{\iconComp[20px]{}}{Spawn \enemyMidnightBeast{} attacking the PCs.}
}

\subsubsection{Allied Help}
\label{ssec:middnightMotherAllies}

When a PC spends \iconOpp{} to call an ally for help, they can choose one of the effects below. Please note, specified ally must be present in the scene to use their ability:

\scriptedOption{\npcShallan{} - Lightweaving}{
	Triggers Fear of Lightweaving of all Midnight creatures. They gain disadvantage on their attack rolls until the end of their turn. 
}

\scriptedOption{\npcAdolin{} - Strike: Shardblade}{
	Attack +10, any one target. Graze: 9 (2d8) spirit damage. Hit: 22 (2d8 + 10) spirit damage. \\
	\textit{(assume \npcAdolin{} always has enough Focus to graze)}
}

\scriptedOption{\npcRenarin{} - Alter Fortune}{
	The PC rolls a d20 and stores the result. As a \iconReaction{} they can replace  with that value one d20 roll made by a character they can see.
}

\scriptedOption{Bridgeman - Defensive Position}{
	The PC can Brace without a shield or a cover.
}

\scriptedOption{\npcJasnah{} - Living Transformation}{
	\npcJasnah{} makes a Transformation attack against a Spiritual Defense of an enemy: \\
	Attack +10, any one target. Graze: 19 (3d12) spirit damage. Hit: 29 (3d12 + 10) spirit damage.\\
	\textit{(assume \npcJasnah{} always has enough Focus to graze)}
}

\scriptedOption{\npcKaladin{} - Gravitational Slam}{	
	\npcKaladin{} makes a Gravitation (+8) test against a Physical Defense of an enemy. On a success, it moves towards another enemy. Each of them takes 5 (1d10) impact damage for each 10 ft. the first enemy moved as a result of this ability (with a minimum of 1d10).
}

\scriptedOption{\npcNavani{}'s Fabrial - Overcharge}{
	You can use a fabrial from \npcNavani{} as \iconFreeAction{}.
}

\scriptedOption{\npcDalinar{}'s Officer - Confident Command}{
	The PC's movement increases by 10 ft. and they ignore Exhausted, Slowed, and Surprised. The PC also gains a d6 die they can use in their next test.
}

\scriptedOption{Ghostblood - Strike: Dueling Cane}{
	Attack +6, any one target. Graze: 3 (1d6) impact damage. Hit: 11 (1d6 + 8) impact damage, and the target is knocked Prone. When the target is knocked Prone, strike them again. \\
	\textit{(assume he always has enough Focus to graze)}
}

\scriptedOption{\enemyDiagramSurgeon{} - First Aid}{
	The PC recovers 10 (1d10 + 5) Health.
}

\scriptedOption{Ardent - Galvanize}{
	The PC rolls their recovery die and recovers that much Focus.
}

\scriptedOption{Son of Honor - Instill Confidence}{
	The PC becomes Focused.
}

\subsection{Rewards}
\label{ssec:middnightMotherRewards}

After defeating \npcMidnightMotherFull{} award each PC with 1 level. 

Depending on the previous interactions, various NPCs may reward the PCs with their patronage or different rewards.

From this point, \fullref{ssec:chapter8runningEvent} is ended and no longer affects the party.  On top of that, all rolls for random encounters in \fullref{sec:darkTunnels} gain advantage and the GM cannot spend \iconComp{} to remove the advantage.

If the PC persuaded \npcMidnightMother{} that they are her friends, award them with a pet - a friendly \enemyMidnightBeast{} will follow them everywhere. If it is slain, it will regenerate after the next Long Rest.




\subsection{Ignored...}
\label{ssec:middnightMotherIgnored}

It is possible for the players to completely ignore the issue of \npcMidnightMother{}. They can allow \npcShallan{} to find the Unmade on her own and allow her to face the threat without the support from the PCs.

If that happens, allow \npcShallan{} to defeat the Unmade and end \fullref{ssec:chapter8runningEvent}, but do not reward the PCs with other rewards. That means they will have 1 fewer levels after finishing this chapter.